\documentclass[12pt,letterpaper]{hmcpset}
\usepackage{amsmath}
\usepackage{amssymb}
\usepackage{amsfonts}
\usepackage{graphicx} % Required for inserting images
\usepackage{hyperref}
\usepackage{enumitem}
\usepackage[margin=1in]{geometry}
\usepackage{graphicx}
\usepackage{float}
\DeclareMathOperator{\vspan}{span}
\newcommand{\bb}[1]{\mathbb{#1}}
\newcommand{\bvec}[1]{\mathbf{#1}}
\newcommand{\pardx}[2]{\frac{\partial #1}{\partial #2}}
\newcommand{\ddx}[2]{\frac{d #1}{d #2}}
\newcommand{\norm}[1]{\Vert #1 \Vert}
\newcommand{\abs}[1]{\vert #1 \vert}
\newcommand{\set}[1]{\{ #1 \}}
\setlist[enumerate, 1]{label = (\alph*)}

% info for header block in upper right hand corner
\name{Jayden Thadani}
\class{Linear Algebra Done Right self-study}
\assignment{Exercises 6.A---Inner Products and Norms}

\begin{document}

\problemlist{}

\begin{problem}[29]
    For \( u , v \in V \), define \( d(u, v) = \Vert u - v \Vert \)
    \begin{enumerate}
        \item Show that $d$ is a metric on $V$.
        \item Show that if $V$ is finite-dimensional, then $d$ is a complete metric on $V$ (meaning that every Cauchy sequence converges).
        \item Show that every finite-dimensional subspace of $V$ is a closed subset of $V$.
    \end{enumerate}
\end{problem}

\begin{solution}
    \begin{enumerate}
        \item In order to show that $d$ is a metric, we need to show that it is positive definite, symmetric, and satisfies the triangle inequality.\\\\
        First we show that $d$ is positive definite. Let \( u, v, w \) be any vectors in $V$. Since \( \Vert w \Vert = 0 \) if and only if \( w = 0 \), we have that \( d (u, v) = 0 \) if and only if \( u - v = 0 \). Thus, we have \( d (u, v) = 0 \) if and only if \( u = v \). Note also that \( \Vert w \Vert \ge 0 \), giving us \( d(u, v) \ge 0 \) always. Combining these two statements about $d$ gives us that it is positive definite.\\\\
        Now we show that $d$ is symmetric. This follows obviously from the properties of the norm of a scaled vector. For any \( u, v \in V \)
        \[
        d(u, v) = \Vert u - v \Vert = \Vert (-1) (v - u) \Vert = \vert -1 \vert \Vert v - u \Vert = \Vert v - u \Vert = d(v, u).
        \]\\
        Finally, we only need to point out an obvious conversion between what the triangle means for norms and what it means for metric spaces. For any \( u, v, w \in V \) 
        \[
        d(u, v) + d(v, w) = \Vert v - u \Vert + \Vert w - v \Vert.
        \]
        By the triangle inequality (for norms) we have
        \[
        \norm{v - u} + \norm{w - v} \le \norm{(v - u) + (w - v)} = \norm{w - u}.
        \]
        Thus, we have the triangle inequality for metrics:
        \[
        d(u, v) + d(v, w) \le d(u, w)
        \]
        \item If $V$ is finite-dimensional, then we have an orthonormal basis \(u_1, \ldots, u_m\). (This is proven later in the book). Let $\set{v_n}$ be any Cauchy sequence in $V$.\\\\
        Then we can write each element of the sequence
        \[
        v_n = a_{n, 1} u_1 + \ldots + a_{n, m} u_m.
        \]\\
        Since $\set{v_n}$ is Cauchy, we have that for any \( \varepsilon > 0 \), there exists some \( N \in \bb{N} \) such that for all \( j, k > N \) we have \( d(v_j, v_k) < \varepsilon \). We can rewrite the expression as
        \[
        d(a_{j, 1} u_1 + \ldots + a_{j, m} u_m, a_{k, 1} u_1 + \ldots + a_{k, m} u_m) = \norm{ (a_{j, 1} - a_{k, 1}) u_1 + \ldots + (a_{j, m} - a_{k, m}) u_m}.
        \]
        Thus,
        \[
        \norm{ (a_{j, 1} - a_{k, 1}) u_1 + \ldots + (a_{j, m} - a_{k, m}) u_m} < \varepsilon
        \]
        (with all the same qualifiers for \( j, k, \varepsilon \) ).
        Since the basis vectors are orthogonal, and the right and left hand sides of the inequality are obviously nonnegative, we get
        \[
         \abs{a_{j, 1} - a_{k, 1}}^2 \norm{u_1}^2 + \ldots + \abs{a_{j, m} - a_{k, m}}^2 \norm{u_m}^2 < \varepsilon^2
        \]
        and since the basis vectors are also unit vectors, we can simplify this to
        \[
        \abs{a_{j, 1} - a_{k, 1}}^2 + \ldots + \abs{a_{j, m} - a_{k, m}}^2 < \varepsilon^2
        \]
        Since all the terms on the left hand side are nonnegative, we can write
        \[
        \abs{a_{j, p} - a{k, p}}^2 < \varepsilon^2
        \]
        for any \( p \in \set{1, \ldots, m} \).
        Again, since both sides of the inequality are positive, we get
        \[
        \abs{a_{j, p} - a_{k, p}} < \varepsilon.
        \]\\
        Specifically, for any \( p \in \set{1, \ldots, m} \), for any \( \varepsilon > 0 \), there exists some \( N \in \bb{N} \) such that for all \( j, k > N \)
        \[
        \abs{a_{j, p} - a_{k, p}} < \varepsilon.
        \]
        That is, for any \( p \in \set{1, \ldots m} \) the sequence \( \set{a_{n, p}} \) is Cauchy.\\\\
        Since $V$ is a field over $\bb{R}$ or $\bb{C}$, both of which are complete, for all $p$, $\set{a_{n, p}}$ converges to some $A_p$. That is, for any \( \varepsilon' > 0 \), there is some \( N \in \bb{N} \) such that for all \( n > N \), \( \abs{A_p - a_{n, p}} < \varepsilon' \). This motivates us to think about the sequence of vectors converging to \( A_1 u_1 + \ldots + A_m u_m \).\\\\
        For any \( \varepsilon' > 0 \) we can have some $N$ such that for all $n > N$, for all $p$, \( \abs{A_p - a_{n, p}} < \frac{\varepsilon}{m}' \) (if for each sequence of coefficients we get $N_p$, then take some $n$ greater than \( \max{N_p} \)). We consider
        \[
        d(A_1 u_1 + \ldots + A_m u_m, a_{n, 1} u_1 + \ldots a_{n, m} u_m) = \norm{(A_1 u_1 + \ldots + A_m u_m) - (a_{n, 1} u_1 + \ldots a_{n, m} u_m)}
        \]
        and perform the same expansion as previously to get
        \[
        d(A_1 u_1 + \ldots + A_m u_m, a_{n, 1} u_1 + \ldots a_{n, m} u_m) \le \abs{A_1 - a_{n, 1}} + \ldots + \abs{A_m - a_{n, m}} < m \frac{\varepsilon'}{m} = \varepsilon'
        \]\\
        Specifically, for any \( \varepsilon' > 0 \) there is some \( N \in \bb{N} \) such that for all \( n > N \)
        \[
        d(A_1 u_1 + \ldots + A_m u_m, a_{n, 1} u_1 + \ldots a_{n, m} u_m) < \varepsilon'.
        \]
        That is, \( \set{v_n} \) converges to \( A_1 u_1 + \ldots A_m u_m \).
        \item By the construction in part (b), any Cauchy sequence in a subspace $U$ of $V$ converges to a point in $U$. Thus, $U$ contains all of its limits points, and thus, is closed.
    \end{enumerate}
\end{solution}

\begin{problem}[30]
	Fix a positive integer $n$. The \textbf{\textit{Laplacian}} $\Delta p$ of a twice differentiable function $p$ on $\bb{R}^n$ is the function on $\bb{R}^n$ defined by
	\[
		\Delta p = \frac{\partial^2 p}{\partial x_1} + \ldots + \frac{\partial^2 p}{\partial x_n}
	\]
A function $p$ is called \textbf{\textit{harmonic}} if \( \Delta p = 0 \).\\\\
	A \textbf{\textit{polynomial}} on $\bb{R}^n$ is a linear combination of functions of the form \( x_1^{m_1} \cdots x_n^{m_n} \), where \( m_1 \ldots m_n \) are nonnegative integers.\\\\
	Suppose $q$ is a polynomial on $\bb{R}^n$. Prove that there exists a harmonic polynomial $p$ on $\bb{R}^n$ such that \( p(x) = q(x) \) for every \( x \in \bb{R}^n \) with \( \Vert x \Vert = 1 \).

\end{problem}

\begin{solution}
	Let $\mathcal{P}(\bb{R}^n)$ denote the space of polynomials on $\bb{R}^n$ and let $\mathcal{P}_m(\bb{R}^n)$ be the space of polynomials on $\bb{R}^n$ of degree at most $m$.
	If $q \in \mathcal{P}(\bb{R}^n)$, and we want some harmonic polynomial $p$ such that \( p = q \) for every \( x \in \bb{R}^n \) with \( \norm{x} = 1 \), we can reasonably think that
	\[
		p = q + (1 - \norm{x}^2)r
	\]
	where $r$ is some polynomial. Such a $p$ would necessarily give us \( p = q \) at all the desired points, but there might not be such a harmonic polynomial. Thus, we need to show that there is some harmonic polynomial of the same form. That is, we need to show that there exists some $r$ such that
	\[
		\Delta p = \Delta q + \Delta ((1 - \norm{x}^2)r) = 0.
	\]
	This amounts to showing that there exists some $r$ such that
	\[
		\Delta ((1 - \norm{x}^2)r) = -\Delta q.
	\]
	This, can be done by considering the linear map giving \( r \mapsto \Delta ((1 - \norm{x}^2)r) \) and showing that it is surjective. However, since the space of all polynomials is infinite-dimensional, this would require a direct proof and would not amount to any simplification. Instead, we can choose the right finite-dimensional space so that the map is an operator and we only have to prove that it is injective.\\\\
	Note that if \( \operatorname{deg} q  = m + 2 \), then \( \operatorname{deg} \Delta q = m \). Similarly, if \( \operatorname{deg} r = m \), then \( \deg (1 - \norm{x}^2) r  = m + 2 \) and thus \( \deg \Delta (1 - \norm{x}^2) r  = m \). Thus, we can just consider the space of polynomials of degree at most $m$.\\\\
	Specifically, we can consider \( T \in \mathcal{L}(\mathcal{P}_m(\bb{R}^n))\) defined by
	\[
		Tr = \Delta (1 - \norm{x}^2) r.
	\]
	To be done, we only have to show that $T$ is injective.\\\\
	Suppose \( Ts = 0 \). Then, using the product rule for the Laplacian we have
	\[
		s \Delta (1 - \norm{x}^2) + (1 - \norm{x}^2) \Delta s + \nabla s \cdot \nabla (1 - \norm{x}^2) = 0.
	\]
	Calculating all the derivatives where we can, we get
	\[
		s (-2n) + (1 - \norm{x}^2) \Delta s + \sum_{j = 1}^n s_{x_j} (-2 x_j) = 0
	\]
	which can be rewritten as
	\[
		s = \frac{1}{2n} (1 - \norm{x}^2) \Delta s - \frac{1}{n} \sum_{j = 1}^n \pardx{s}{x_j} x_j
	\]
\end{solution}

\begin{problem}[31]
	Use inner products to prove Apollonius's identity: in a triangle with sides of length $a$, $b$, and $c$, let $x$ be the length of the line segment from the midpoint of the side of length $c$ to the opposite vertex. Then 
	\[
		a^2 + b^2 = \frac{1}{2} c^2 + d^2
	\]
\end{problem}

\begin{solution}
	Consider the obvious choice of $v_a$, $v_b$, $v_c$, and $v_d$ such that \( \norm{v_\alpha} = \alpha \). Note that this then gives us
	\[
		\frac{v_c}{2} + v_d = v_a = v_c + v_b
	\]
	and
	\[
		\frac{v_c}{2} + v_b = v_d.
	\]
	Thus, we have 
	\[
		\norm{\frac{v_c}{2} + v_d}^2 = \norm{v_a}^2
	\]
	which we can expand using the dot product to get
	\[
		\frac{c^2}{4} + d^2 + \frac{v_c}{2} \cdot v_d = a^2.
	\]
	We also have 
	\[
		\norm{\frac{v_c}{2} + v_b}^2 = \norm{v_d}^2
	\]
	which gives us
	\[
		\frac{c^2}{4} + b^2 + \frac{v_c}{2} \cdot v_b = d^2.
	\]
	Swapping sides and adding the equations gives us
	\[
		2d^2 + \frac{v_c}{2} \cdot v_d = a^2 + b^2 + \frac{v_c}{2} \cdot v_b
	\]
	which we can rewrite as 
	\[
		2d^2 + \frac{v_c}{2} \cdot (v_d - v_b) = a^2 + b^2
	\]
	\[
		2d^2 + \frac{v_c}{2} \cdot v_c = a^2 + b^2
	\]
	which finally gives us
	\[
		a^2 + b^2 = 2d^2 + \frac{c^2}{2}
	\]
	as desired.
\end{solution}

% Add pairs of problems and solutions as needed

\end{document}
